\section{Methods}\label{sec:methods}

\subsection{Study design and reporting}
This is a secondary analysis of de-identified, publicly released data. We report the prediction
analysis according to TRIPOD+AI~\citep{collins2024tripodai} and the descriptive analysis according to
STROBE~\citep{vonelm2007strobe}; the completed TRIPOD+AI checklist is in Supplementary
Table~\ref{tab:tripod}. Figure~\ref{fig:framework} shows the study framework in three lanes: data
preparation (A), model development and validation (B), and the checks that test the results (C).
Before running any experiment, we fixed the feature sets, models, metrics, expected directions of
effect and re-weighting targets in a configuration file (random seed \Seed). All tables, figures
and numbers in this paper are generated directly from the results files by code.

\begin{figure*}[t]\centering
\resizebox{\linewidth}{!}{% Figure 1 -- study framework (three lanes). Counts are macros from generated/numbers.tex.
\begin{tikzpicture}[
  font=\sffamily\scriptsize,
  box/.style={draw=#1!65!black, fill=#1!7, rounded corners=2pt, align=left, inner sep=3pt,
              text width=0.28\linewidth, minimum height=24mm},
  lane/.style={font=\sffamily\scriptsize\bfseries, text=#1!60!black, rotate=90, anchor=south},
  arr/.style={-{Latex[length=1.8mm]}, thick, draw=black!50},
  bus/.style={thick, draw=black!50},
]
\matrix (m) [row sep=8mm, column sep=5mm, ampersand replacement=\&] {
\node[box=blue] (src) {\textbf{1 Open sources} (no approval, no DUA)\\
  cBioPortal: METABRIC, MSK-2018, MBCproject, Sanger, Broad (Mexico/Vietnam), BCCRC, TCGA\\
  TCIA: Duke-Breast-Cancer-MRI\\
  UCI: Ljubljana\\
  \emph{SHA-256 manifest; no redistribution}}; \&
\node[box=blue] (cdm) {\textbf{2 Common data model} (one row/patient)\\
  Outcome: invasive Nottingham grade 1--3\\
  In-situ disease excluded, never relabelled\\
  ILC kept separate from IDC; men excluded\\
  Quarantined: NPI, component scores, stage}; \&
\node[box=blue] (fs) {\textbf{3 Analysis sets}\\
  \NGradedTotal{} graded patients, \NGradedCohorts{} cohorts\\
  Full: age, menopause, ER, PR, HER2, histology, T, N\\
  Core: age, ER, PR, HER2, histology}; \\
\node[box=orange] (e1) {\textbf{E1 Describe} (STROBE)\\
  Kendall $\tau_b$, bias-corrected Cram\'er's $V$, Spearman $\rho$\\
  Bootstrap 95\% CIs; Benjamini--Hochberg\\
  Proportional-odds model}; \&
\node[box=orange] (e2) {\textbf{E2 Develop} in METABRIC ($n$=\NMetabric)\\
  Nested \Repeats$\times$\Folds-fold CV; inner log-loss tuning\\
  Baseline, LR, ordinal LR, $k$-NN, RF, XGBoost, LightGBM\\
  In-fold impute$\to$scale$\to$model\\
  Multiple imputation ($m$=\MIm); complete case}; \&
\node[box=orange] (e3) {\textbf{E3/E5 Validate} (TRIPOD+AI)\\
  Frozen models: MSK-2018, Duke (full); MBCproject, Sanger, Broad (core)\\
  Leave-one-dataset-out\\
  Calibration, intercept update, decision curves}; \\
\node[box=teal] (e4) {\textbf{E4 Leakage audit}\\
  NPI as predictor\\ SMOTE/imputation before split\\ Component scores (Sanger, Duke)\\ Record-level split}; \&
\node[box=teal] (e6) {\textbf{E6/E9 Plausibility and imaging}\\
  TreeSHAP directions vs pre-specified expectations; ordinal odds ratios\\
  Duke: clinical vs DCE-MRI radiomics vs fusion}; \&
\node[box=violet] (e7) {\textbf{E7 Transport to Saudi case-mix}\\
  Rake cohorts to published Saudi ER, PR, HER2, histology, age\\
  Implied vs published grade; pooled anchors; age sensitivity}; \\
};
\node[draw=black!55, fill=black!4, rounded corners=2pt, anchor=north, inner sep=3pt, align=center,
      text width=0.92\linewidth] (out) at ([yshift=-6mm]m.south)
  {\textbf{Output:} \texttt{make all} $\to$ \texttt{results/} $\to$ every table, figure and in-text number (\texttt{numbers.tex})};
\draw[arr] (src) -- (cdm); \draw[arr] (cdm) -- (fs);
% bus from data lane to analysis lane
\coordinate (b1) at ($(fs.south)!0.5!(e3.north)$);
\draw[bus] (fs.south) -- (b1) -- (b1 -| e1.north);
\foreach \n in {e1,e2,e3} \draw[arr] (b1 -| \n.north) -- (\n.north);
\draw[arr] (e2) -- (e3);
% bus from analysis lane to audit lane
\coordinate (b2) at ($(e2.south)!0.5!(e6.north)$);
\draw[bus] (e2.south) -- (b2); \draw[bus] (b2 -| e4.north) -- (b2 -| e7.north);
\draw[bus] (e3.south) -- (b2 -| e3.south);
\foreach \n in {e4,e6,e7} \draw[arr] (b2 -| \n.north) -- (\n.north);
\draw[arr] (m.south) -- (out.north);
\node[lane=blue]   at ([xshift=-1mm]src.west) {Data};
\node[lane=orange] at ([xshift=-1mm]e1.west)  {Analysis};
\node[lane=teal]   at ([xshift=-1mm]e4.west)  {Audits \& transport};
\end{tikzpicture}
}
\caption{Study framework. \textbf{A}, open data are downloaded with checksums and mapped to one row
per patient under four harmonisation rules; variables that contain grade are set aside.
\textbf{B}, associations are measured with statistics suited to each variable type; models are
developed in METABRIC with nested cross-validation, then frozen and validated in five external
cohorts and across datasets. \textbf{C}, the results are checked for leakage and clinical
plausibility, extended with MRI radiomics in the Duke cohort, and transported to published Saudi
case-mix. E1--E10 are the experiment stages in the code repository.}\label{fig:framework}
\end{figure*}

\subsection{Data sources}
All data were downloaded without registration from the cBioPortal data
hub~\citep{cerami2012cbio,gao2013cbio}, The Cancer Imaging Archive (TCIA) and the UCI Machine Learning
Repository. A download script records a SHA-256 checksum for every file, so any later change by a data
custodian can be detected. The cohorts and their roles were:
\begin{itemize}[leftmargin=*,itemsep=1pt,topsep=2pt]
\item \textbf{METABRIC} (UK/Canada; development)~\citep{curtis2012metabric,pereira2016metabric}: a
surgical cohort with Nottingham grade, tumour size in millimetres, number of positive nodes,
receptor status and menopausal status.
\item \textbf{MSK-2018} (USA; external validation, full feature set)~\citep{razavi2018msk}: patients
sequenced at Memorial Sloan Kettering, enriched for advanced hormone-receptor-positive disease, with
grade, T and N category, receptors and menopausal status at diagnosis.
\item \textbf{Duke-Breast-Cancer-MRI} (USA; external validation, full feature set; imaging
extension)~\citep{saha2018duke}: 922 patients with invasive cancer imaged at Duke University
Hospital between 2000 and 2014, with Nottingham grade and its component scores, receptors, T and N
category, menopausal status and 529 pre-computed MRI features (CC BY-NC 4.0; obtained from a
checksummed public copy of the TCIA files).
\item \textbf{MBCproject} (USA/Canada; external, core features)~\citep{mbcproject2022,jain2023mbcproject}:
a patient-partnered cohort of metastatic disease with diagnosis-time records abstracted from charts;
age is given in five-year bands.
\item \textbf{Sanger} (UK/Europe; external, core features)~\citep{stephens2012sanger}: 100 primary
tumours with grade and the three Nottingham component scores.
\item \textbf{Broad} (Mexico and Vietnam; external, core features)~\citep{banerji2012broad}: the only
open grade-labelled cohort from low- and middle-income countries that we found.
\item \textbf{BCCRC} (Canada; triple-negative; descriptive only)~\citep{shah2012bccrc} and
\textbf{TCGA-BRCA} (USA; descriptive only, no structured grade)~\citep{tcga2012brca}.
\item \textbf{Ljubljana} (Slovenia; sensitivity analysis)~\citep{zwitter1988ljubljana}: \NLj{}
patients with a ``degree of malignancy'' of 1--3 and grouped covariates.
\end{itemize}
We found no openly available, grade-labelled patient-level cohort from Saudi Arabia or the Gulf.
Saudi data therefore enter the study only as published summary figures
(Section~\ref{sec:transport}).

\subsection{Harmonisation to a common data model}
We mapped every source to one row per patient with the following variables: age in years,
post-menopausal status, ER, PR and HER2 status (positive or negative; equivocal results treated as
missing), histological type (ductal/no special type, lobular, mixed ductal--lobular, other special
type), T category (T1 $\le$20\,mm, T2 21--50\,mm, T3 or T4) and nodal status (any positive regional
node, including micrometastasis, versus none). Four rules prevented common data-handling errors:
(i) the outcome is invasive Nottingham grade 1--3, and any other value is treated as missing;
(ii) pure in-situ disease (DCIS, LCIS) and non-carcinomas are excluded, never relabelled as invasive;
(iii) lobular carcinoma is never merged with ductal carcinoma; and (iv) men are excluded. When a
patient had several sequenced samples (MSK-2018), we kept the diagnosis-time record once.
Supplementary Table~\ref{tab:flow} shows the participant flow.

\subsection{Outcome, predictors and feature sets}
The primary outcome was grade in three classes; the main secondary outcome was grade~3 versus
grades 1--2. We used two pre-specified feature sets. The \emph{full} set (age, menopausal status,
ER, PR, HER2, histology, T category and nodal status) was available in METABRIC, MSK-2018 and Duke.
The \emph{core} set (age, ER, PR, HER2 and histology) was available in all external cohorts.
Variables that contain grade by construction---the Nottingham Prognostic Index and the three
Nottingham component scores---and stage were never used as predictors; they appear only in the
leakage audit. A unit test stops the analysis if any of them enters a feature set.

\subsection{Descriptive and association statistics}
We matched each association measure to the type of variable: Kendall's $\tau_b$ for grade against
binary or ordinal variables, bias-corrected Cram\'er's $V$~\citep{bergsma2013cramer} for histological
type, and Spearman's $\rho$ with the Kruskal--Wallis $\varepsilon^2$ for age. Every estimate has a
bootstrap 95\% confidence interval (\Boot{} resamples), and $p$ values were adjusted for the
\NTests{} tests with the Benjamini--Hochberg procedure~\citep{benjamini1995fdr}. These associations
describe which tumours tend to be high grade; they do not imply causation. A proportional-odds
(ordinal logistic) model fitted to complete cases in METABRIC summarises the associations jointly.

\subsection{Model development}
We compared seven models of increasing flexibility: a baseline that predicts the class prevalences,
multinomial logistic regression (LR) with an $L_2$ penalty, proportional-odds logistic regression,
$k$-nearest neighbours, random forest, XGBoost~\citep{chen2016xgboost} and
LightGBM~\citep{ke2017lightgbm}. We also pre-specified a tabular foundation model,
TabPFN~\citep{hollmann2025tabpfn}, as an optional eighth model (status: \TabpfnStatus; see
Limitations). Imputation of missing values~\citep{vanbuuren2011mice}, standardisation and the model
itself were combined in one pipeline, so every step was fitted on training data only. We tuned
hyper-parameters by grid search, minimising log-loss in an inner \InnerFolds-fold cross-validation.
We did not rebalance the classes, because resampling distorts predicted
probabilities~\citep{vandengoorbergh2022imbalance}. Instead, when a single predicted class was
needed, we chose the class with the largest ratio of predicted probability to training prevalence,
$\hat y=\arg\max_k p_k/\pi_k$.

\paragraph{Missing data} The primary analysis used single chained-equation imputation fitted within
each training fold. We ran two pre-specified sensitivity analyses: multiple imputation with
$m=\MIm$ imputed datasets, a separate model for each and the average of their predicted
probabilities~\citep{sisk2023imputation,hoogland2020missing}; and a complete-case analysis. Both
used the same resampling design as the primary analysis.

\paragraph{Sample size} We calculated the minimum development sample size with the criteria of Riley
et al.~\citep{riley2020samplesize}, applied to each pair of grades as recommended for multinomial
models~\citep{pate2023multinomial}. The criteria require a shrinkage factor of at least 0.9, a small
difference between apparent and adjusted $R^2$ ($\delta$=0.05) for every pairwise sub-model with
\SSParams{} predictor parameters, and precise estimation of each grade's prevalence ($\pm$0.05). The
expected pairwise $R^2$ values came from out-of-fold LR predictions, which are not optimistic.

\subsection{Validation}
Internal validation used nested, stratified \Folds-fold cross-validation in METABRIC, repeated
\Repeats{} times. We averaged each patient's out-of-fold predictions across repeats and computed
bootstrap 95\% confidence intervals (\Boot{} resamples) on these averages. Each model was compared
with LR by a paired bootstrap of the difference in macro-AUC, using the same resamples for both
models. For external validation, each model was tuned on all of METABRIC, frozen, and applied
unchanged to MSK-2018 and Duke (full set) and to MSK-2018, Duke, MBCproject, Sanger and Broad (core
set). We also used leave-one-dataset-out validation: models were trained on all but one of METABRIC,
MSK-2018, Duke, MBCproject and Sanger and tested on the remaining
cohort~\citep{steyerberg2016validation}.

\subsection{Performance measures}
We measured discrimination with the macro-averaged one-versus-rest AUC and the AUC for each grade;
classification with macro-F1, balanced accuracy and quadratic-weighted $\kappa$ (QWK); and the overall
accuracy of the predicted probabilities with the multiclass Brier score. For grade~3 versus 1--2 we
also report the calibration intercept (whether predictions are too high or too low on average), the
calibration slope (whether they are too extreme or too modest), smoothed calibration
curves~\citep{vancalster2019calibration} and decision curves~\citep{vickers2006dca}. For all three
grades we compare the observed and mean predicted prevalence~\citep{vanhoorde2014multinomial}. In
each external cohort we tested a simple model update: re-estimating only the intercept, cross-fitted
over two halves of the cohort. This separates miscalibration caused by a different grade prevalence
from miscalibration of the predictor effects~\citep{riley2024evalpart2}.

\subsection{Leakage audit}
We compared the correct protocol with four common shortcuts on the same patients and
models~\citep{kapoor2023leakage}: (a) adding the Nottingham Prognostic Index, which includes grade,
as a predictor; (b) imputing and oversampling with SMOTE~\citep{chawla2002smote} on the whole dataset
before cross-validation; (c) adding the Nottingham component scores as predictors (Sanger and Duke);
and (d) splitting by record instead of by patient when some patients have several records
(MSK-2018). The last shortcut is the tabular version of splitting image tiles rather than patients in
digital pathology~\citep{bussola2019slipping}.

\subsection{Interpretability and clinical plausibility}
We computed TreeSHAP values for the grade-3 output of XGBoost~\citep{lundberg2017shap,lundberg2020tree}.
Before the analysis we wrote down the expected direction of effect for seven features: ER-negative,
PR-negative, HER2-positive, larger tumours, node-positive disease, younger age and non-lobular
histology should all be associated with higher grade. A direction was counted as confirmed when the
bootstrap interval of the Spearman correlation between a feature's value and its SHAP value excluded
zero in the expected direction.

\subsection{Imaging extension: MRI radiomics in Duke}
In the graded Duke patients we compared three penalised multinomial LR models using nested
cross-validation (\Repeats{} repeats of \Folds{} folds): the clinical full set, the 529 pre-computed
dynamic contrast-enhanced MRI features~\citep{saha2018duke}, and both combined. The penalty was tuned
in an inner 3-fold loop. The value added by radiomics was the paired-bootstrap difference in
macro-AUC compared with the clinical model.

\subsection{Transport to Saudi case-mix}\label{sec:transport}
Without Saudi patient data, we asked a narrower question: what grade distribution, and what model
predictions, would the Western cohorts show if their patients had the receptor, histology and age
profile reported for Saudi patients? We re-weighted the METABRIC, MSK-2018 and Duke patients by
iterative proportional fitting (raking) so that their weighted profile matched published Saudi
figures~\citep{steingrimsson2023transport}. The primary target was a series of Saudi nationals
treated at King Fahd University Hospital, Dammam~\citep{kussaibi2021saudi}, which reports ER, PR,
HER2 and histology. Because it reports only the median age (51 years), we assumed that 50\% of
patients were aged $\le$50 years and varied this from \AgeSensLo\% to \AgeSensHi\% in a sensitivity
analysis. The second target was a tertiary-centre series from King Fahad Medical City,
Riyadh~\citep{durrani2025healthcare}, which reports ER, PR, HER2, histology and the share aged
$\le$50 years; we recalculated its receptor percentages from the published counts. Every value was
transcribed from the published tables, with table and page numbers recorded in the repository.
For each target and cohort we report the re-weighted grade distribution with a bootstrap 95\%
confidence interval (\TransportBoot{} resamples, each re-weighted afresh), the effective sample size
of the weights, and the weighted mean predicted probabilities. We pooled the published Saudi grade-3
proportions with a DerSimonian--Laird random-effects model on the logit scale, using every series that
reports grade for invasive cancer only~\citep{kussaibi2021saudi,durrani2025healthcare,albasri2014madinah};
a sensitivity analysis also included a small series that grades all lesions~\citep{alamri2024najran}.
This analysis shows how much of the difference in grade can be explained by measured case-mix; it
cannot capture local grading practice or unmeasured differences.

\subsection{Software and reproducibility}
We used Python 3.13 with scikit-learn 1.9.1, imbalanced-learn 0.14.2, XGBoost 3.4.1, LightGBM 4.7.0,
SHAP 0.54.1 and statsmodels 0.15.0 (versions are pinned in the repository). A single command
(\texttt{make all}) downloads and checks the data, runs the unit tests that enforce the harmonisation
and leakage rules, runs experiments E1--E10 and regenerates every table, figure and number in the
paper.
